% GENERATED FILE. Edit canonical lessons, then run npm run generate:books.
% canonical-source-hash: fnv1a64:22299546af97f051
% canonical-lessons: KA-C307-godamu, KA-C307-phalaka, KA-C307-sayi, KA-C307-simesunna, KA-C307-skelu, KA-R307-first-pass-words-from-chapters-299-302, KA-R307-second-pass-words-from-chapters-303-307

\chapter{Warehouse, Signboard, Ink, Chalk, Ruler}
\label{ch:ka-warehouse-signboard-ink-chalk-ruler}
\hlchaptermodality{307}

\begin{chapteropening}
\textbf{By the end of this chapter:} \emph{I can say a warehouse, a signboard, ink, chalk and a ruler.}

Complete the last lesson of chapter 307: I can say a warehouse, a signboard, ink, chalk and a ruler.
\end{chapteropening}

\section[\texorpdfstring{\kn{ಗೋದಾಮು} (gōdāmu)}{ಗೋದಾಮು (gōdāmu)}]{\kn{ಗೋದಾಮು} (gōdāmu) — a warehouse, a godown}

\label{lesson:KA-C307-godamu}



\begin{quote}
Before the new one: say the Kannada for a plumber, then the Kannada for a gardener.
\end{quote}

\subsection*{You'll want to know: \kn{ಗೋದಾಮು}}
\textbf{\kn{ಗೋದಾಮು}} — \emph{gōdāmu} — \textquotedblleft{}a warehouse, a godown\textquotedblright{}.

\begin{grammarlens}[title={What you've built}]
One more word for this chapter.
\end{grammarlens}

\subsection*{Your turn}
\begin{itemize}
\raggedright
  \item \emph{Say it:} \emph{gōdāmu}
  \item \emph{Say it:} \emph{gōdāmu}, once more
  \item \emph{Say it:} say \emph{gōdāmu}
  \item \emph{Recall:} say the Kannada for a plumber, then the Kannada for a gardener, then say \emph{gōdāmu} again
\end{itemize}

\begin{tcolorbox}[breakable,colback=teal!4,colframe=teal!35!black,title={Before you move on}]
What does \kn{ಗೋದಾಮು} mean? (\textquotedblleft{}A warehouse, a godown\textquotedblright{}.) Say it once more.
\end{tcolorbox}
\section[\texorpdfstring{\kn{ಫಲಕ} (phalaka)}{ಫಲಕ (phalaka)}]{\kn{ಫಲಕ} (phalaka) — a signboard, a nameplate}

\label{lesson:KA-C307-phalaka}



\begin{quote}
Before the new one: say the Kannada for a gardener, then the Kannada for a warehouse.
\end{quote}

\subsection*{You'll want to know: \kn{ಫಲಕ}}
\textbf{\kn{ಫಲಕ}} — \emph{phalaka} — \textquotedblleft{}a signboard, a nameplate\textquotedblright{}.

\begin{grammarlens}[title={What you've built}]
One more word for this chapter.
\end{grammarlens}

\subsection*{Your turn}
\begin{itemize}
\raggedright
  \item \emph{Say it:} \emph{phalaka}
  \item \emph{Say it:} \emph{phalaka}, once more
  \item \emph{Say it:} say \emph{phalaka}
  \item \emph{Recall:} say the Kannada for a gardener, then the Kannada for a warehouse, then say \emph{phalaka} again
\end{itemize}

\begin{tcolorbox}[breakable,colback=teal!4,colframe=teal!35!black,title={Before you move on}]
What does \kn{ಫಲಕ} mean? (\textquotedblleft{}A signboard, a nameplate\textquotedblright{}.) Say it once more.
\end{tcolorbox}
\section[\texorpdfstring{\kn{ಶಾಯಿ} (śāyi)}{ಶಾಯಿ (śāyi)}]{\kn{ಶಾಯಿ} (śāyi) — ink}

\label{lesson:KA-C307-sayi}



\begin{quote}
Before the new one: say the Kannada for a warehouse, then the Kannada for a signboard.
\end{quote}

\subsection*{You'll want to know: \kn{ಶಾಯಿ}}
\textbf{\kn{ಶಾಯಿ}} — \emph{śāyi} — \textquotedblleft{}ink\textquotedblright{}.

\begin{grammarlens}[title={What you've built}]
One more word for this chapter.
\end{grammarlens}

\subsection*{Your turn}
\begin{itemize}
\raggedright
  \item \emph{Say it:} \emph{śāyi}
  \item \emph{Say it:} \emph{śāyi}, once more
  \item \emph{Say it:} say \emph{śāyi}
  \item \emph{Recall:} say the Kannada for a warehouse, then the Kannada for a signboard, then say \emph{śāyi} again
\end{itemize}

\begin{tcolorbox}[breakable,colback=teal!4,colframe=teal!35!black,title={Before you move on}]
What does \kn{ಶಾಯಿ} mean? (\textquotedblleft{}Ink\textquotedblright{}.) Say it once more.
\end{tcolorbox}
\section[\texorpdfstring{\kn{ಸೀಮೆಸುಣ್ಣ} (sīmesuṇṇa)}{ಸೀಮೆಸುಣ್ಣ (sīmesuṇṇa)}]{\kn{ಸೀಮೆಸುಣ್ಣ} (sīmesuṇṇa) — chalk}

\label{lesson:KA-C307-simesunna}



\begin{quote}
Before the new one: say the Kannada for a signboard, then the Kannada for ink.
\end{quote}

\subsection*{You'll want to know: \kn{ಸೀಮೆಸುಣ್ಣ}}
\textbf{\kn{ಸೀಮೆಸುಣ್ಣ}} — \emph{sīmesuṇṇa} — \textquotedblleft{}chalk\textquotedblright{}.

\begin{grammarlens}[title={What you've built}]
One more word for this chapter.
\end{grammarlens}

\subsection*{Your turn}
\begin{itemize}
\raggedright
  \item \emph{Say it:} \emph{sīmesuṇṇa}
  \item \emph{Say it:} \emph{sīmesuṇṇa}, once more
  \item \emph{Say it:} say \emph{sīmesuṇṇa}
  \item \emph{Recall:} say the Kannada for a signboard, then the Kannada for ink, then say \emph{sīmesuṇṇa} again
\end{itemize}

\begin{tcolorbox}[breakable,colback=teal!4,colframe=teal!35!black,title={Before you move on}]
What does \kn{ಸೀಮೆಸುಣ್ಣ} mean? (\textquotedblleft{}Chalk\textquotedblright{}.) Say it once more.
\end{tcolorbox}
\section[\texorpdfstring{\kn{ಸ್ಕೇಲು} (skēlu)}{ಸ್ಕೇಲು (skēlu)}]{\kn{ಸ್ಕೇಲು} (skēlu) — a ruler (for drawing lines)}

\label{lesson:KA-C307-skelu}



\begin{quote}
Before the new one: say the Kannada for ink, then the Kannada for chalk.
\end{quote}

\subsection*{You'll want to know: \kn{ಸ್ಕೇಲು}}
\textbf{\kn{ಸ್ಕೇಲು}} — \emph{skēlu} — \textquotedblleft{}a ruler (for drawing lines)\textquotedblright{}.

\begin{grammarlens}[title={What you've built}]
One more word for this chapter.
\end{grammarlens}

\subsection*{Your turn}
\begin{itemize}
\raggedright
  \item \emph{Say it:} \emph{skēlu}
  \item \emph{Say it:} \emph{skēlu}, once more
  \item \emph{Say it:} say \emph{skēlu}
  \item \emph{Recall:} say the Kannada for ink, then the Kannada for chalk, then say \emph{skēlu} again
\end{itemize}

\begin{tcolorbox}[breakable,colback=teal!4,colframe=teal!35!black,title={Before you move on}]
What does \kn{ಸ್ಕೇಲು} mean? (\textquotedblleft{}A ruler (for drawing lines)\textquotedblright{}.) Say it once more.
\end{tcolorbox}
\section[(dialogue)]{Words from chapters 299-302 — a first pass}

\label{lesson:KA-R307-first-pass-words-from-chapters-299-302}



\begin{quote}
Say the words for chalk and a ruler.
\end{quote}

\subsection*{Your turn}
Say these aloud:

\begin{itemize}
\raggedright
  \item an old woman, people, a human being, a stranger, elders
  \item a photo, a camera, a programme, a TV serial, a novel
  \item glass, plastic, a brick, a list, a griddle
  \item brass, firewood, wool, cement, bamboo
\end{itemize}

\begin{tcolorbox}[breakable,colback=teal!4,colframe=teal!35!black,title={Before you move on}]
An old woman? (\textbf{\kn{ಮುದುಕಿ}}, \emph{muduki}.) People? (\textbf{\kn{ಜನ}}, \emph{jana}.) A human being? (\textbf{\kn{ಮನುಷ್ಯ}}, \emph{manuṣya}.) A stranger? (\textbf{\kn{ಅಪರಿಚಿತ}}, \emph{aparicita}.) Elders? (\textbf{\kn{ಹಿರಿಯರು}}, \emph{hiriyaru}.) A photo? (\textbf{\kn{ಫೋಟೋ}}, \emph{phōṭō}.) A camera? (\textbf{\kn{ಕ್ಯಾಮರಾ}}, \emph{kyāmarā}.) A programme? (\textbf{\kn{ಕಾರ್ಯಕ್ರಮ}}, \emph{kāryakrama}.) A TV serial? (\textbf{\kn{ಧಾರಾವಾಹಿ}}, \emph{dhārāvāhi}.) A novel? (\textbf{\kn{ಕಾದಂಬರಿ}}, \emph{kādambari}.) Glass? (\textbf{\kn{ಗಾಜು}}, \emph{gāju}.) Plastic? (\textbf{\kn{ಪ್ಲಾಸ್ಟಿಕ್}}, \emph{plāsṭik}.) A brick? (\textbf{\kn{ಇಟ್ಟಿಗೆ}}, \emph{iṭṭige}.) A list? (\textbf{\kn{ಪಟ್ಟಿ}}, \emph{paṭṭi}.) A griddle? (\textbf{\kn{ಕಾವಲಿ}}, \emph{kāvali}.) Brass? (\textbf{\kn{ಹಿತ್ತಾಳೆ}}, \emph{hittāḷe}.) Firewood? (\textbf{\kn{ಕಟ್ಟಿಗೆ}}, \emph{kaṭṭige}.) Wool? (\textbf{\kn{ಉಣ್ಣೆ}}, \emph{uṇṇe}.) Cement? (\textbf{\kn{ಸಿಮೆಂಟ್}}, \emph{simeṇṭ}.) Bamboo? (\textbf{\kn{ಬಿದಿರು}}, \emph{bidiru}.)
\end{tcolorbox}
\section[(dialogue)]{Words from chapters 303-307 — a second pass}

\label{lesson:KA-R307-second-pass-words-from-chapters-303-307}



\begin{quote}
Say the words for a raised veranda and a loft.
\end{quote}

\subsection*{Your turn}
Say these aloud:

\begin{itemize}
\raggedright
  \item a raised veranda, a loft, a door latch, a handle, a smell
  \item a lesson, an exam, training, a subject, mathematics
  \item a temple priest, an artist, a singer, an actor, a journalist
  \item a cobbler, a potter, a mason, a plumber, a gardener
  \item a warehouse, a signboard, ink, chalk, a ruler
\end{itemize}

\begin{tcolorbox}[breakable,colback=teal!4,colframe=teal!35!black,title={Before you move on}]
A raised veranda? (\textbf{\kn{ಜಗಲಿ}}, \emph{jagali}.) A loft? (\textbf{\kn{ಅಟ್ಟ}}, \emph{aṭṭa}.) A door latch? (\textbf{\kn{ಚಿಲಕ}}, \emph{cilaka}.) A handle? (\textbf{\kn{ಹಿಡಿಕೆ}}, \emph{hiḍike}.) A smell? (\textbf{\kn{ವಾಸನೆ}}, \emph{vāsane}.) A lesson? (\textbf{\kn{ಪಾಠ}}, \emph{pāṭha}.) An exam? (\textbf{\kn{ಪರೀಕ್ಷೆ}}, \emph{parīkṣe}.) Training? (\textbf{\kn{ತರಬೇತಿ}}, \emph{tarabēti}.) A subject? (\textbf{\kn{ವಿಷಯ}}, \emph{viṣaya}.) Mathematics? (\textbf{\kn{ಗಣಿತ}}, \emph{gaṇita}.) A temple priest? (\textbf{\kn{ಅರ್ಚಕ}}, \emph{arcaka}.) An artist? (\textbf{\kn{ಕಲಾವಿದ}}, \emph{kalāvida}.) A singer? (\textbf{\kn{ಗಾಯಕ}}, \emph{gāyaka}.) An actor? (\textbf{\kn{ನಟ}}, \emph{naṭa}.) A journalist? (\textbf{\kn{ಪತ್ರಕರ್ತ}}, \emph{patrakarta}.) A cobbler? (\textbf{\kn{ಮೋಚಿ}}, \emph{mōci}.) A potter? (\textbf{\kn{ಕುಂಬಾರ}}, \emph{kumbāra}.) A mason? (\textbf{\kn{ಮೇಸ್ತ್ರಿ}}, \emph{mēstri}.) A plumber? (\textbf{\kn{ಪ್ಲಂಬರ್}}, \emph{plambar}.) A gardener? (\textbf{\kn{ತೋಟಗಾರ}}, \emph{tōṭagāra}.) A warehouse? (\textbf{\kn{ಗೋದಾಮು}}, \emph{gōdāmu}.) A signboard? (\textbf{\kn{ಫಲಕ}}, \emph{phalaka}.) Ink? (\textbf{\kn{ಶಾಯಿ}}, \emph{śāyi}.) Chalk? (\textbf{\kn{ಸೀಮೆಸುಣ್ಣ}}, \emph{sīmesuṇṇa}.) A ruler? (\textbf{\kn{ಸ್ಕೇಲು}}, \emph{skēlu}.)
\end{tcolorbox}
